% Chương 1 (sửa trực tiếp; bản tiếng Việt)
\chapter{GIỚI THIỆU}

\section{Động lực}

Một giáo viên đứng trước 30-50 học sinh không thể theo dõi sự chú ý của từng em vào mọi thời điểm, và một chỉ số đơn giản ở cấp lớp cho biết có bao nhiêu học sinh đang có vẻ mất tập trung có thể giúp giáo viên điều chỉnh nhịp độ bài giảng. Các hệ thống dùng camera làm việc này đã tồn tại, nhưng những hệ thống nổi tiếng nhất, ở các trường học Trung Quốc giai đoạn 2018-2019, nhận diện khuôn mặt của mọi học sinh, gửi video lên máy chủ và chấm điểm từng em cho giáo viên và phụ huynh. Chúng gây ra làn sóng phản đối của dư luận và một hệ thống thậm chí đã bị cơ quan chức năng đình chỉ (chương 2).

Câu hỏi luận văn đặt ra là liệu có thể xây dựng một chỉ số tổng hợp về sự chú ý mà không cần giám sát hay không. Các lựa chọn thiết kế xuất phát từ đó:

\begin{itemize}
\tightlist
\item
  \textbf{Tư thế thay cho khuôn mặt.} Chỉ trích xuất các điểm khóa của cơ thể (13 điểm phần thân trên); không lưu khuôn mặt nào, và học sinh chỉ là các số chỗ ngồi ẩn danh trong một phiên.
\item
  \textbf{Biên thay cho đám mây.} Mọi xử lý chạy trên một Raspberry Pi 5 với bộ tăng tốc Hailo-8 đặt ngay trong lớp. Không khung hình nào rời khỏi thiết bị, và mặc định không khung hình nào được lưu.
\item
  \textbf{Cấp lớp thay cho điểm số cá nhân.} Đầu ra dành cho giáo viên là tỷ lệ học sinh trong lớp có vẻ mất tập trung, kèm một khoảng bất định, chứ không phải báo cáo theo từng học sinh.
\end{itemize}

Các vấn đề còn lại mang tính kỹ thuật: mô hình tư thế gọn nhẹ bỏ sót các học sinh nhỏ ở hàng cuối, và việc ``nhìn đi chỗ khác'' khó đánh giá từ bộ khung xương khi lớp học và vị trí camera khác nhau.

\section{Câu hỏi nghiên cứu}

\begin{enumerate}
\def\labelenumi{\arabic{enumi}.}
\tightlist
\item
  \textbf{RQ1, tính khả thi trên thiết bị biên.} Một quy trình đánh giá mức độ tập trung dựa trên tư thế có thể chạy thời gian thực cho cả lớp học trên Raspberry Pi 5 với Hailo-8 hay không, và chi phí về điện năng và nhiệt là bao nhiêu?
\item
  \textbf{RQ2, độ bao phủ.} Một mô hình tư thế gọn nhẹ tìm được bao nhiêu học sinh trong một lớp học thực, và tiền xử lý cùng hậu xử lý có thể nâng con số đó mà không phải liên tục huấn luyện lại mô hình hay không?
\item
  \textbf{RQ3, quyết định cho từng học sinh.} Các quyết định dựa trên tư thế, bằng một bộ luật đã hiệu chỉnh và bằng các bộ phân loại huấn luyện trên một bộ dữ liệu lớp học công khai, phân biệt học sinh tập trung với học sinh mất tập trung trong lớp học thực tốt đến đâu?
\item
  \textbf{RQ4, chỉ số cấp lớp.} Có thể xây dựng một tỷ lệ cấp lớp với khoảng bất định đáng tin cậy khi nhiều học sinh bị bỏ sót hay không?
\end{enumerate}

\section{Đóng góp}

\begin{itemize}
\tightlist
\item
  \textbf{hcpose}, một runtime mã nguồn mở cho các mô hình YOLOv8-pose tùy biến trên Hailo-8, gồm một engine GStreamer và một engine Python cho kết quả như nhau, cài đặt và chạy bằng một lệnh.\footnote{Mã nguồn và các bản phát hành: \url{https://github.com/LTH3ar/hailo-custom-yolo-pose-runtime}}
\item
  \textbf{Chia ô ảnh cho các hàng xa:} chạy suy luận bổ sung trên các ô phóng to ở cuối lớp, gộp kết quả không trùng lặp bằng C++ ngay trong pipeline GStreamer. Trên ảnh SCB, kỹ thuật này nâng tỷ lệ học sinh được phát hiện từ 55\% lên 72\%, và từ 25\% lên 66\% ở một phần ba xa nhất, mà không thay đổi mô hình.
\item
  \textbf{Loại bỏ khung xương trùng lặp:} một bước hậu xử lý loại bỏ hộp thứ hai mà YOLO-pose thường sinh ra cho cùng một học sinh.
\item
  \textbf{Bộ phân tích mức độ tập trung} với hiệu chỉnh theo từng chỗ ngồi và hai thước đo góc quay đầu. Thước đo thứ hai, độ lệch của mũi so với trung điểm hai vai, vẫn hoạt động khi nhìn nghiêng, trường hợp mà cách ước lượng thông thường dựa trên tai thất bại.
\item
  \textbf{Ước lượng cấp lớp} có ghi nhớ chỗ ngồi, gán trọng số cho các vùng gần, giữa và xa, và báo cáo khoảng Wilson, với cách đếm chỗ ngồi chống lại các phát hiện giả.
\item
  \textbf{Đánh giá trên bốn loại bằng chứng:}

  \begin{itemize}
  \tightlist
  \item
    hai bộ dữ liệu lớp học công khai, SCB và ClaBehavior
  \item
    bốn đoạn video lớp học thực (ba đoạn công khai, một đoạn ghi hình cho luận văn này)
  \item
    451 nhãn mù do người gán từ ba đoạn video
  \item
    một buổi học dàn dựng theo kịch bản
  \end{itemize}

  Kết hợp lại, chúng cho thấy các bộ phân loại huấn luyện trên các bộ dữ liệu này không chuyển giao tốt sang các buổi giảng thực, kể cả khi có lớp ngẩng đầu chú ý, trong khi hiệu chỉnh theo từng học sinh thì có.
\item
  \textbf{Đo đạc trên thiết bị biên:} thông lượng, độ trễ, CPU, nhiệt độ và công suất của bo mạch.
\item
  \textbf{Phân tích quyền riêng tư} đối chiếu với Luật Bảo vệ dữ liệu cá nhân của Việt Nam (2025) và Nghị định 356/2025.
\end{itemize}

\section{Phạm vi}

Hệ thống ước lượng sự chú ý có thể quan sát được từ tư thế và hướng đầu; nó không đo lường việc học, cảm xúc hay mức độ tham gia nhận thức, và không thể thấy một chiếc điện thoại được cầm như cầm bút. Các mô hình tư thế đã được huấn luyện trước luận văn này và giữ cố định: mọi cải tiến nằm ở runtime, hậu xử lý và phân tích. Phần đánh giá dùng các đoạn ghi hình ngắn và các buổi học dàn dựng thay vì một đợt triển khai thực địa dài.

\section{Cấu trúc luận văn}

Chương 2 tổng quan về phát hiện mức độ tập trung, ước lượng tư thế trên thiết bị biên và các hệ thống đã triển khai từng gây lo ngại về quyền riêng tư. Chương 3 mô tả hệ thống. Chương 4 đánh giá hệ thống và phân tích các tính chất về quyền riêng tư và tính pháp lý. Chương 5 trả lời các câu hỏi nghiên cứu và nêu hướng phát triển.
